\section{Introduction}

South Africa faces a severe early-grade reading crisis where 81\% of Grade 4 learners tested in PIRLS 2021 could not read for meaning in any language \citep{spaull2023pirls}. Within the Schooling 2030 agenda, DBE-led initiatives have sought to improve early-grade literacy through structured pedagogy programmes combining lesson plans, teacher training and integrated reading materials \citep{dbe_egrs}. These instructional supports address one component of learners’ educational environment. Another potentially important feature is learners’ exposure to bullying by their peers. According to the PIRLS 2021 results, approximately three quarters of South African learners reported experiencing bullying monthly or weekly \citep{mullis2023pirls}. The coexistence of poor reading achievement and widespread bullying raises an important educational and economic question about whether learners’ exposure to bullying is associated with the development of their foundational reading skills.
 Existing evidence generally finds a negative relationship between exposure to bullying or school violence and educational achievement \citep{anton_erxleben2016bullying, ozyildirim2024effect}. Using PIRLS 2021 data, this essay examines the conditional association between bullying frequency and Grade 4 reading achievement in South Africa and Brazil. Brazil is used as a comparison country to assess whether the relationship observed in South Africa also appears in another unequal upper-middle-income economy with a comparable level of GDP per capita but substantially higher reading achievement \citep{de2023comparative, spaull2023pirls}. The comparison therefore provides a benchmark for assessing whether the South African association is context-specific, although it does not hold all country-level differences constant. The central research question is: Is bullying frequency associated with Grade 4 reading achievement in South Africa, does this association differ by gender and home socioeconomic status, and how does the South African pattern compare with Brazil?
 Five nested regression specifications are estimated separately for each country, beginning with bullying frequency and progressively adding learner, household and school characteristics. Gender and home SES interactions are then used to examine whether the association differs across learners, while common-sample and alternative-measure specifications assess the robustness of the results. 
